\chapter{Mellékletek}
A mellékletek a WaccApp frontendjének gyakorlati megvalósítását kiegészítő anyagokat tartalmazzák. A fejezetben szereplő mappaszerkezet, API-végpont-kivonat, manuális tesztforgatókönyvek, kérdőív-adatszerkezeti példa és képernyőképes nézetbemutatások célja az, hogy a dolgozat törzsszövegében bemutatott frontendműködés konkrét projektállományokhoz és ellenőrizhető felületi megoldásokhoz kapcsolódjon.

\section{A WaccApp frontend mappaszerkezete}
A melléklet a WaccApp projekt frontend szempontból lényeges állományait foglalja össze. A teljes forráskód a digitális mellékletben szerepel.
\begin{CodeBlock}
webhook-main/
|-- app.js
|-- questionnaire.js
|-- whatsapp_messages.db
|-- package.json
|-- kerdoivek/
|   |-- olvasas.json
|   `-- orvosi.json
`-- public/
    |-- index.html
    |-- conv.html
    |-- filter.html
    |-- messages.html
    |-- sent-messages.html
    |-- questionnaire.html
    |-- template.html
    |-- template-menu.html
    |-- template-single.html
    |-- login.html
    |-- menu.html
    |-- admin.html
    |-- account.html
    |-- privacy.html
    |-- send.html
    |-- questionnaires.json
    |-- uploads/
    `-- sent_media/
\end{CodeBlock}

\section{Frontend szempontból fontos API-végpontok kivonata}

A következő táblázat a WaccApp fontosabb API-végpontokat foglalja össze. A cél bemutatni a fő frontendnézetek milyen kapcsolódási pontokon keresztül kommunikálnak a háttérrendszerrel.

\begin{table}[H]
\centering
\caption{Frontend szempontból fontos API-végpontok kivonata.}
\label{tab:melleklet-api-vegpontok}
\footnotesize
\setlength{\tabcolsep}{3pt}
\renewcommand{\arraystretch}{1.25}

\begin{tabularx}{\textwidth}{|>{\raggedright\arraybackslash}p{0.23\textwidth}|>{\raggedright\arraybackslash}p{0.12\textwidth}|>{\raggedright\arraybackslash}p{0.25\textwidth}|Y|}
\hline
\textbf{Végpont} &
\textbf{Módszer} &
\textbf{Kapcsolódó frontend nézet} &
\textbf{Frontend szerep} \\
\hline

\path{/send-message} &
POST &
\path{index.html} &
Egyedi szöveges üzenet küldése. \\
\hline

\path{/send-template} &
POST &
\path{index.html}, \path{template.html} &
Sablonos üzenet küldése. \\
\hline

\path{/send-file-message} &
POST &
\path{index.html} &
Média- vagy fájlcsatolmány küldése. \\
\hline

\path{/schedule} &
POST &
\path{index.html} &
Időzített szöveges, sablonos vagy kérdőíves küldés. \\
\hline

\path{/schedule-media} &
POST &
\path{index.html} &
Időzített médiaküldés. \\
\hline

\path{/messages} &
GET &
\path{messages.html}, \path{conv.html} &
Beérkező vagy tárolt üzenetek lekérése. \\
\hline

\path{/sent-messages} &
GET &
\path{sent-messages.html}, \path{conv.html} &
Elküldött üzenetek lekérése. \\
\hline

\path{/contacts} &
GET &
\path{conv.html}, \path{filter.html} &
Kontaktok, név--telefonszám kapcsolat kezelése. \\
\hline

\path{/all-questionnaires} &
GET &
\path{index.html}, \path{questionnaire.html} &
Kérdőívek listázása. \\
\hline

\path{/save-to-questionnaire} &
POST &
\path{questionnaire.html}, \path{template.html} &
Kérdőívlogika mentése. \\
\hline

\end{tabularx}
\end{table}

\section{Manuális tesztforgatókönyvek}

A mellékletben szereplő tesztforgatókönyvek a főszövegben bemutatott tesztelési szempontrendszer részletesebb bontását adják. A cél az volt, hogy a WaccApp legfontosabb frontendfolyamatai konkrét bemenetekkel, lépésekkel, elvárt eredményekkel és tényleges tapasztalatokkal is ellenőrizhetők legyenek.

\clearpage
\begingroup
\pagestyle{fancy}
\begin{landscape}
\scriptsize
\setlength{\tabcolsep}{2pt}
\renewcommand{\arraystretch}{1.18}
\begin{longtable}{|
>{\raggedright\arraybackslash}p{0.050\linewidth}|
>{\raggedright\arraybackslash}p{0.125\linewidth}|
>{\raggedright\arraybackslash}p{0.145\linewidth}|
>{\raggedright\arraybackslash}p{0.160\linewidth}|
>{\raggedright\arraybackslash}p{0.170\linewidth}|
>{\raggedright\arraybackslash}p{0.170\linewidth}|
>{\raggedright\arraybackslash}p{0.055\linewidth}|}
\caption{Részletes manuális tesztforgatókönyvek.}
\label{tab:reszletes-manualis-tesztek}\\
\hline
\textbf{Azon.} &
\textbf{Tesztelt folyamat} &
\textbf{Kiindulási állapot} &
\textbf{Bemenet vagy művelet} &
\textbf{Elvárt eredmény} &
\textbf{Tényleges eredmény} &
\textbf{Státusz} \\
\hline
\endfirsthead

\multicolumn{7}{c}{\tablename\ \thetable{} -- folytatás az előző oldalról} \\
\hline
\textbf{ID} &
\textbf{Tesztelt folyamat} &
\textbf{Kiindulási állapot} &
\textbf{Bemenet vagy művelet} &
\textbf{Elvárt eredmény} &
\textbf{Tényleges eredmény} &
\textbf{Státusz} \\
\hline
\endhead

\hline
\multicolumn{7}{r}{Folytatás a következő oldalon} \\
\endfoot

\hline
\endlastfoot

T1 &
Egyedi szöveges üzenet küldése &
Az \path{index.html} oldal elérhető, a backend fut &
Érvényes telefonszám és nem üres üzenetszöveg megadása, majd küldés indítása &
A rendszer elküldi az üzenetet, sikeres visszajelzést ad, és az üzenet később visszakereshető &
A küldés lefutott, a felület sikeres állapotot jelzett, az üzenet megjelent az elküldött tételek között &
Megfelelt \\
\hline

T2 &
Üres üzenet kezelése &
Az \path{index.html} oldalon címzett megadható &
Telefonszám megadása után az üzenetmező üresen hagyása &
A frontend ne indítsa el a küldést, hanem jelezze a hiányzó üzenettartalmat &
A küldés nem indult el, a felület hibajelzést adott &
Megfelelt \\
\hline

T3 &
Hibás telefonszám kezelése &
Az \path{index.html} oldalon a küldési űrlap elérhető &
Hibás formátumú telefonszám megadása &
A frontend jelezze, hogy a telefonszám nem megfelelő formátumú &
A felület hibát jelzett, a küldés nem indult el &
Megfelelt \\
\hline

T4 &
Sablonos üzenet küldése &
Elérhető legalább egy sablon &
Sablon kiválasztása, szükséges paraméterek kitöltése, küldés indítása &
A rendszer a kiválasztott sablon alapján elindítja a küldést, majd visszajelzést ad &
A sablonos küldés elindult, a felület visszajelzést adott &
Megfelelt \\
\hline

T5 &
Hiányzó sablonparaméter kezelése &
Sablonküldési felület megnyitva &
Sablon kiválasztása, de kötelező paraméter üresen hagyása &
A frontend jelezze, hogy a sablonos küldéshez hiányzik adat &
A küldés nem indult el, a felület hibás állapotot jelenített meg &
Megfelelt \\
\hline

T6 &
Kérdőív létrehozása &
A \path{questionnaire.html} oldal elérhető &
Kérdés, válaszlehetőségek és következő lépések megadása &
A kérdőív menthető legyen, és később indítható formában rendelkezésre álljon &
A kérdőív szerkezete menthető volt, a folyamat később kiválasztható maradt &
Megfelelt \\
\hline

T7 &
Válaszfüggő kérdőívútvonal ellenőrzése &
Létező kérdőív indítható &
Eltérő válaszok kipróbálása egy kérdőívfolyamatban &
A megadott válasz alapján a megfelelő következő kérdés következzen &
A kérdőív a megadott válaszlogika szerint haladt tovább &
Megfelelt \\
\hline

T8 &
Beérkező üzenetek listázása &
Léteznek beérkezett üzenetek &
\path{messages.html} megnyitása &
A beérkező üzenetek listaalapú formában jelenjenek meg &
A beérkezett kommunikációs események megjelentek a nézetben &
Megfelelt \\
\hline

T9 &
Elküldött üzenetek listázása &
Léteznek elküldött üzenetek &
\path{sent-messages.html} megnyitása &
Az elküldött tételek időponttal és tartalommal áttekinthetők legyenek &
Az elküldött üzenetek listában megjelentek &
Megfelelt \\
\hline

T10 &
Beszélgetésnézet ellenőrzése &
Léteznek beérkező és elküldött üzenetek ugyanahhoz a kontakthoz &
\path{conv.html} megnyitása adott kontakt alapján &
A rendszer időrendi, chat-szerű formában jelenítse meg a kommunikációt &
A beérkező és kimenő üzenetek elkülönítve, beszélgetésként jelentek meg &
Megfelelt \\
\hline

T11 &
Médiatartalom megjelenítése beszélgetésben &
Korábban elküldött kép vagy dokumentum rendelkezésre áll &
\path{conv.html} megnyitása az adott kommunikációs előzménnyel &
A kép vizuálisan, a dokumentum megnyitható hivatkozásként jelenjen meg &
A médiatartalom a beszélgetési kontextusban értelmezhetően megjelent &
Megfelelt \\
\hline

T12 &
Név és telefonszám szinkronizálása szűrésnél &
A \path{filter.html} oldalon vannak ismert kontaktok &
Név kiválasztása a szűrőmezőben &
A kapcsolódó telefonszám automatikusan igazodjon a kiválasztott névhez &
A név és telefonszám mezők szinkronban működtek &
Megfelelt \\
\hline

T13 &
Dátumintervallum szerinti szűrés &
Léteznek különböző időpontú kommunikációs események &
Kezdő- és záródátum megadása &
Csak az adott időszakba eső találatok jelenjenek meg &
A találati lista az időszak szerint szűkült &
Megfelelt \\
\hline

T14 &
Kérdőíves válaszösszesítés &
Létezik kérdőíves kommunikáció és válaszadat &
Kérdőív és kérdés kiválasztása a \path{filter.html} oldalon &
A felület darabszám szerint jelenítse meg az ismert válaszokat és az „Egyéb” kategóriát &
A válaszok összesítése megjelent, a nem ismert válaszok külön kategóriába kerültek &
Megfelelt \\
\hline

T15 &
Időzített küldés rögzítése &
Az \path{index.html} oldalon az időzítési mező elérhető &
Érvényes címzett, tartalom és jövőbeli időpont megadása &
A rendszer rögzítse az időzített műveletet, és erről visszajelzést adjon &
Az időzített művelet rögzítése után a felület visszajelzést adott &
Megfelelt \\
\hline

T16 &
Hibás időzítési adat kezelése &
Időzített küldési folyamat indítható &
Hiányzó vagy múltbeli időpont megadása &
A frontend jelezze, hogy az időzítés nem megfelelő &
A művelet nem indult el, a felület javítandó hibát jelzett &
Megfelelt \\
\hline

T17 &
Reszponzív megjelenés ellenőrzése &
A fő nézetek elérhetők &
\path{index.html}, \path{conv.html}, \path{filter.html} és kérdőívszerkesztő oldalak megnyitása kisebb kijelzőn &
A fő műveletek kezelhetők, az üzenetek olvashatók, az űrlapok követhetők maradjanak &
A fő nézetek használhatók maradtak, kisebb finomítási lehetőségekkel &
Megfelelt \\
\hline

T18 &
Regressziós ellenőrzés módosítás után &
Egy korábbi funkció módosítása megtörtént &
Kapcsolódó funkciók újbóli kipróbálása, például kérdőívindítás, beszélgetésnézet és szűrés &
A módosítás ne rontsa el a kapcsolódó működéseket &
A kapcsolódó folyamatok továbbra is használhatók maradtak &
Megfelelt \\

\end{longtable}
\end{landscape}
\endgroup
\clearpage
\pagestyle{fancy}

A részletes tesztforgatókönyvek alapján látható, hogy az ellenőrzés nemcsak az egyes oldalak megnyitására irányult, hanem a WaccApp fő frontendfolyamatainak gyakorlati kipróbálására is. A tesztek között szerepeltek sikeres műveletek, hibás adatbeviteli helyzetek, visszakeresési folyamatok, kérdőíves útvonalak, médiamegjelenítés, időzítés és reszponzív használat is. Ez azért fontos, mert a rendszer használhatósága azon múlik, hogy ezek a folyamatok együtt is követhető és javítható működést adnak.

\section{Példa kérdőív-adatszerkezet}

A WaccApp kérdőíves működésének egyik fontos sajátossága, hogy a kérdések nem egyszerű lineáris sorrendben követik egymást, hanem a válaszokhoz kapcsolódó következő lépések alapján. Az alábbi egyszerűsített példa azt mutatja be, hogyan jelenhet meg egy válaszfüggő kérdőívlogika adatszerkezeti formában.

\begin{CodeBlock}
const questionnaires = {
    kiszallitas: {
        kerdoiv1: {
            next: {
                'Igen': 'kerdoiv2',
                'Nem': 'kerdoiv2'
            }
        },
        kerdoiv2: {
            next: {
                'Igen': 'kerdoiv3',
                'Nem': 'kerdoiv3'
            }
        },
        kerdoiv3: {
            next: {
                'Megtartom': null,
                'Nem tartom meg': 'csomag_gond'
            }
        },
        csomag_gond: {
            next: null
        }
    }
};
\end{CodeBlock}

A példában a kiszallitas nevű kérdőív több egymáshoz kapcsolódó lépésből áll. A next objektum azt határozza meg, hogy egy adott válasz után melyik következő kérdés vagy állapot következzen. A null érték a folyamat lezárását jelzi. A teljes kérdőívlogika a digitális mellékletben található forráskódban szerepel.